\subsection{Selección de diseño conceptual}\label{SS:CritSel}

Mediante el Proceso de Análisis Jerárquico (AHP), el cual es un método de selección multicriterio se procedió a evaluar las distintas propuestas de solución que buscan resolver todas las necesidades identificadas previamente.

Para poder realizar la selección de los diferentes componentes se consideró la siguiente serie de criterios:

\begin{itemize}
    \item Costo:
    Mientras el precio sea menor se le asigna una mayor calificación.
    \item Facilidad de manufactura:
    Se refiere a que tan fácil es manufacturar, en caso de ser necesario, el componente a utilizar.
    \item Disponibilidad:
    Es la facilidad con la que se puede conseguir algún componente en el mercado.
    \item Seguridad de uso:
    Representa el riesgo que puede experimentar para una persona con determinado sistema o componente, mientras más peligroso menor ponderación obtiene.
    \item Facilidad de uso:
    Mientras más sencillo sea de utilizar por personas sin los conocimientos técnicos, mejor evaluación.
    \item Integración:
    Hace referencia a la facilidad de interconectar un componente con otro, se busca tener un alto nivel de integración entre componentes.
    \item Durabilidad:
    Se refiere al tiempo de vida útil de cada componente.
    \item Mantenimiento:
    Facilidad con la cual pueda realizarse el mantenimiento de cierto componente, a menor dificultad mayor calificación.
    \item Ergonomía:
    Se trata de la relación y comodidad que se genera al utilizar algún componente, se busca una alta ergonomía.
    \item Portabilidad:
    Facilidad de desplazar un componente de un lugar a otro, se busca una alta portabilidad.
    \item Estética
    Agrado visual que genera cierto componente, sin tener en cuenta su funcionabilidad.
\end{itemize}

Estos diferentes criterios fueron ponderados, utilizando el criterio AHP, en una tabla, como se muestra en el apéndice \ref{Ap:Crit}, para determinar qué características son las más sobresalientes y que tendrán mayor peso al momento de realizar la selección de los elementos del sistema mecatrónico.